% =====================================================================
%  Anatomía de un LLM — Fase 0 — Módulo 0.1
%  Lógica y técnicas de demostración
%  Compilar: latexmk -r ../latexmkrc P01-demostraciones.tex
% =====================================================================
\documentclass[11pt,a4paper,twoside,openany]{book}
\usepackage{../estilo-apuntes}

\newcommand{\ApuntesCurso}{Anatomía de un LLM}
\newcommand{\ApuntesModulo}{Módulo 0.1}

% Notación propia del módulo
\newcommand{\N}{\mathbb{N}}
\newcommand{\Z}{\mathbb{Z}}
\newcommand{\Q}{\mathbb{Q}}
\renewcommand{\V}{\mathrm{V}}% valor de verdad «verdadero» (se redefine \V de la plantilla)
\newcommand{\F}{\mathrm{F}}
\newcommand{\xor}{\oplus}
\newcommand{\divides}{\mid}
\newcommand{\nodivide}{\nmid}
\newcommand{\eqlog}{\equiv}
\DeclarePairedDelimiter{\floor}{\lfloor}{\rfloor}
\DeclarePairedDelimiter{\ceil}{\lceil}{\rceil}

\hypersetup{pdftitle={Módulo 0.1. Lógica y técnicas de demostración (resumen)},pdfauthor={Anatomía de un LLM}}

\makeatletter
\newlabel{ej:proposiciones}{{1.2}{1}{}{ejemplo.2}{}}
\newlabel{ej:proposiciones@cref}{{[ejemplo][2][]1.2}{[1][1][]1}}
\newlabel{ej:reciproco}{{1.7}{2}{}{ejemplo.7}{}}
\newlabel{ej:reciproco@cref}{{[ejemplo][7][]1.7}{[1][2][]2}}
\newlabel{ej:tercero-excluido}{{1.11}{3}{}{ejemplo.11}{}}
\newlabel{ej:tercero-excluido@cref}{{[ejemplo][11][]1.11}{[1][2][]3}}
\newlabel{ej:simplificar}{{1.14}{3}{}{ejemplo.14}{}}
\newlabel{ej:simplificar@cref}{{[ejemplo][14][]1.14}{[1][3][]3}}
\newlabel{obs:1-errores}{{1.17}{4}{Errores frecuentes}{observacion.17}{}}
\newlabel{obs:1-errores@cref}{{[observacion][17][]1.17}{[1][4][]4}}
\newlabel{eje:1-proposiciones}{{1.1}{5}{}{ejercicio.1.1}{}}
\newlabel{eje:1-proposiciones@cref}{{[ejercicio][1][1]1.1}{[1][4][]5}}
\newlabel{eje:1-traducir}{{1.2}{5}{}{ejercicio.1.2}{}}
\newlabel{eje:1-traducir@cref}{{[ejercicio][2][1]1.2}{[1][5][]5}}
\newlabel{eje:1-reciproco}{{1.3}{5}{}{ejercicio.1.3}{}}
\newlabel{eje:1-reciproco@cref}{{[ejercicio][3][1]1.3}{[1][5][]5}}
\newlabel{eje:1-necesaria}{{1.4}{5}{}{ejercicio.1.4}{}}
\newlabel{eje:1-necesaria@cref}{{[ejercicio][4][1]1.4}{[1][5][]5}}
\newlabel{eje:1-silogismo}{{1.5}{5}{}{ejercicio.1.5}{}}
\newlabel{eje:1-silogismo@cref}{{[ejercicio][5][1]1.5}{[1][5][]5}}
\newlabel{eje:1-xor}{{1.6}{5}{}{ejercicio.1.6}{}}
\newlabel{eje:1-xor@cref}{{[ejercicio][6][1]1.6}{[1][5][]5}}
\newlabel{eje:1-equivalencias}{{1.7}{5}{}{ejercicio.1.7}{}}
\newlabel{eje:1-equivalencias@cref}{{[ejercicio][7][1]1.7}{[1][5][]5}}
\newlabel{eje:1-algebra}{{1.8}{5}{}{ejercicio.1.8}{}}
\newlabel{eje:1-algebra@cref}{{[ejercicio][8][1]1.8}{[1][5][]5}}
\newlabel{eje:1-validez}{{1.9}{5}{}{ejercicio.1.9}{}}
\newlabel{eje:1-validez@cref}{{[ejercicio][9][1]1.9}{[1][5][]5}}
\newlabel{eje:1-contar}{{1.10}{5}{}{ejercicio.1.10}{}}
\newlabel{eje:1-contar@cref}{{[ejercicio][10][1]1.10}{[1][5][]5}}
\newlabel{eje:1-leyes}{{1.11}{6}{}{ejercicio.1.11}{}}
\newlabel{eje:1-leyes@cref}{{[ejercicio][11][1]1.11}{[1][6][]6}}
\newlabel{eje:1-demorgan-n}{{1.12}{6}{}{ejercicio.1.12}{}}
\newlabel{eje:1-demorgan-n@cref}{{[ejercicio][12][1]1.12}{[1][6][]6}}
\newlabel{eje:1-completitud}{{1.13}{6}{}{ejercicio.1.13}{}}
\newlabel{eje:1-completitud@cref}{{[ejercicio][13][1]1.13}{[1][6][]6}}
\newlabel{eje:1-bicond-asoc}{{1.14}{6}{}{ejercicio.1.14}{}}
\newlabel{eje:1-bicond-asoc@cref}{{[ejercicio][14][1]1.14}{[1][6][]6}}
\newlabel{ej:traduccion}{{2.4}{7}{}{ejemplo.4}{}}
\newlabel{ej:traduccion@cref}{{[ejemplo][4][]2.4}{[1][7][]7}}
\newlabel{ej:orden}{{2.6}{8}{}{ejemplo.6}{}}
\newlabel{ej:orden@cref}{{[ejemplo][6][]2.6}{[1][8][]8}}
\newlabel{ej:limite}{{2.8}{8}{Convergencia de una sucesión}{ejemplo.8}{}}
\newlabel{ej:limite@cref}{{[ejemplo][8][]2.8}{[1][8][]8}}
\newlabel{ej:negar-limite}{{2.10}{9}{}{ejemplo.10}{}}
\newlabel{ej:negar-limite@cref}{{[ejemplo][10][]2.10}{[1][9][]9}}
\newlabel{obs:2-errores}{{2.12}{9}{Errores frecuentes}{observacion.12}{}}
\newlabel{obs:2-errores@cref}{{[observacion][12][]2.12}{[1][9][]9}}
\newlabel{eje:2-valores}{{2.1}{10}{}{ejercicio.2.1}{}}
\newlabel{eje:2-valores@cref}{{[ejercicio][1][2]2.1}{[1][10][]10}}
\newlabel{eje:2-simbolos}{{2.2}{10}{}{ejercicio.2.2}{}}
\newlabel{eje:2-simbolos@cref}{{[ejercicio][2][2]2.2}{[1][10][]10}}
\newlabel{eje:2-lenguaje}{{2.3}{10}{}{ejercicio.2.3}{}}
\newlabel{eje:2-lenguaje@cref}{{[ejercicio][3][2]2.3}{[1][10][]10}}
\newlabel{eje:2-negar}{{2.4}{10}{}{ejercicio.2.4}{}}
\newlabel{eje:2-negar@cref}{{[ejercicio][4][2]2.4}{[1][10][]10}}
\newlabel{eje:2-continuidad}{{2.5}{10}{}{ejercicio.2.5}{}}
\newlabel{eje:2-continuidad@cref}{{[ejercicio][5][2]2.5}{[1][10][]10}}
\newlabel{eje:2-orden}{{2.6}{10}{}{ejercicio.2.6}{}}
\newlabel{eje:2-orden@cref}{{[ejercicio][6][2]2.6}{[1][10][]10}}
\newlabel{eje:2-distribuir}{{2.7}{10}{}{ejercicio.2.7}{}}
\newlabel{eje:2-distribuir@cref}{{[ejercicio][7][2]2.7}{[1][10][]10}}
\newlabel{eje:2-distribuir-dem}{{2.8}{10}{}{ejercicio.2.8}{}}
\newlabel{eje:2-distribuir-dem@cref}{{[ejercicio][8][2]2.8}{[1][10][]10}}
\newlabel{eje:2-vacio}{{2.9}{10}{}{ejercicio.2.9}{}}
\newlabel{eje:2-vacio@cref}{{[ejercicio][9][2]2.9}{[1][10][]10}}
\newlabel{eje:2-unicidad}{{2.10}{10}{}{ejercicio.2.10}{}}
\newlabel{eje:2-unicidad@cref}{{[ejercicio][10][2]2.10}{[1][10][]10}}
\newlabel{eje:2-acotada}{{2.11}{11}{}{ejercicio.2.11}{}}
\newlabel{eje:2-acotada@cref}{{[ejercicio][11][2]2.11}{[1][11][]11}}
\newlabel{eje:2-arquimedes}{{2.12}{11}{}{ejercicio.2.12}{}}
\newlabel{eje:2-arquimedes@cref}{{[ejercicio][12][2]2.12}{[1][11][]11}}
\newlabel{eje:2-critica}{{2.13}{11}{}{ejercicio.2.13}{}}
\newlabel{eje:2-critica@cref}{{[ejercicio][13][2]2.13}{[1][11][]11}}
\newlabel{eje:2-finito}{{2.14}{11}{}{ejercicio.2.14}{}}
\newlabel{eje:2-finito@cref}{{[ejercicio][14][2]2.14}{[1][11][]11}}
\newlabel{ej:amgm}{{3.7}{13}{Desigualdad con números reales}{ejemplo.7}{}}
\newlabel{ej:amgm@cref}{{[ejemplo][7][]3.7}{[1][13][]13}}
\newlabel{ej:contra-real}{{3.10}{14}{}{ejemplo.10}{}}
\newlabel{ej:contra-real@cref}{{[ejemplo][10][]3.10}{[1][13][]14}}
\newlabel{ej:resto-potencia}{{3.13}{14}{}{ejemplo.13}{}}
\newlabel{ej:resto-potencia@cref}{{[ejemplo][13][]3.13}{[1][14][]14}}
\newlabel{obs:3-errores}{{3.16}{14}{Errores frecuentes}{observacion.16}{}}
\newlabel{obs:3-errores@cref}{{[observacion][16][]3.16}{[1][14][]14}}
\newlabel{eje:3-hipotesis}{{3.1}{15}{}{ejercicio.3.1}{}}
\newlabel{eje:3-hipotesis@cref}{{[ejercicio][1][3]3.1}{[1][15][]15}}
\newlabel{eje:3-esquema}{{3.2}{15}{}{ejercicio.3.2}{}}
\newlabel{eje:3-esquema@cref}{{[ejercicio][2][3]3.2}{[1][15][]15}}
\newlabel{eje:3-cubo}{{3.3}{15}{}{ejercicio.3.3}{}}
\newlabel{eje:3-cubo@cref}{{[ejercicio][3][3]3.3}{[1][15][]15}}
\newlabel{eje:3-polinomio}{{3.4}{15}{}{ejercicio.3.4}{}}
\newlabel{eje:3-polinomio@cref}{{[ejercicio][4][3]3.4}{[1][15][]15}}
\newlabel{eje:3-cuadrados}{{3.5}{15}{}{ejercicio.3.5}{}}
\newlabel{eje:3-cuadrados@cref}{{[ejercicio][5][3]3.5}{[1][15][]15}}
\newlabel{eje:3-producto}{{3.6}{15}{}{ejercicio.3.6}{}}
\newlabel{eje:3-producto@cref}{{[ejercicio][6][3]3.6}{[1][15][]15}}
\newlabel{eje:3-desigualdad}{{3.7}{15}{}{ejercicio.3.7}{}}
\newlabel{eje:3-desigualdad@cref}{{[ejercicio][7][3]3.7}{[1][15][]15}}
\newlabel{eje:3-diferencia}{{3.8}{15}{}{ejercicio.3.8}{}}
\newlabel{eje:3-diferencia@cref}{{[ejercicio][8][3]3.8}{[1][15][]15}}
\newlabel{eje:3-contra-cubo}{{3.9}{16}{}{ejercicio.3.9}{}}
\newlabel{eje:3-contra-cubo@cref}{{[ejercicio][9][3]3.9}{[1][15][]16}}
\newlabel{eje:3-pitagoras}{{3.10}{16}{}{ejercicio.3.10}{}}
\newlabel{eje:3-pitagoras@cref}{{[ejercicio][10][3]3.10}{[1][16][]16}}
\newlabel{eje:3-equivalencia}{{3.11}{16}{}{ejercicio.3.11}{}}
\newlabel{eje:3-equivalencia@cref}{{[ejercicio][11][3]3.11}{[1][16][]16}}
\newlabel{eje:3-real}{{3.12}{16}{}{ejercicio.3.12}{}}
\newlabel{eje:3-real@cref}{{[ejercicio][12][3]3.12}{[1][16][]16}}
\newlabel{eje:3-bicond}{{3.13}{16}{}{ejercicio.3.13}{}}
\newlabel{eje:3-bicond@cref}{{[ejercicio][13][3]3.13}{[1][16][]16}}
\newlabel{eje:3-restos}{{3.14}{16}{}{ejercicio.3.14}{}}
\newlabel{eje:3-restos@cref}{{[ejercicio][14][3]3.14}{[1][16][]16}}
\newlabel{eje:3-division}{{3.15}{16}{}{ejercicio.3.15}{}}
\newlabel{eje:3-division@cref}{{[ejercicio][15][3]3.15}{[1][16][]16}}
\newlabel{ej:no-mayor}{{4.2}{17}{}{ejemplo.2}{}}
\newlabel{ej:no-mayor@cref}{{[ejemplo][2][]4.2}{[1][17][]17}}
\newlabel{ej:no-solucion}{{4.3}{17}{}{ejemplo.3}{}}
\newlabel{ej:no-solucion@cref}{{[ejemplo][3][]4.3}{[1][17][]17}}
\newlabel{obs:4-errores}{{4.12}{19}{Errores frecuentes}{observacion.12}{}}
\newlabel{obs:4-errores@cref}{{[observacion][12][]4.12}{[1][19][]19}}
\newlabel{eje:4-suponer}{{4.1}{19}{}{ejercicio.4.1}{}}
\newlabel{eje:4-suponer@cref}{{[ejercicio][1][4]4.1}{[1][19][]19}}
\newlabel{eje:4-raiz3}{{4.2}{19}{}{ejercicio.4.2}{}}
\newlabel{eje:4-raiz3@cref}{{[ejercicio][2][4]4.2}{[1][19][]19}}
\newlabel{eje:4-raiz-cubica}{{4.3}{19}{}{ejercicio.4.3}{}}
\newlabel{eje:4-raiz-cubica@cref}{{[ejercicio][3][4]4.3}{[1][19][]19}}
\newlabel{eje:4-log}{{4.4}{20}{}{ejercicio.4.4}{}}
\newlabel{eje:4-log@cref}{{[ejercicio][4][4]4.4}{[1][19][]20}}
\newlabel{eje:4-racional-irracional}{{4.5}{20}{}{ejercicio.4.5}{}}
\newlabel{eje:4-racional-irracional@cref}{{[ejercicio][5][4]4.5}{[1][20][]20}}
\newlabel{eje:4-diofantica}{{4.6}{20}{}{ejercicio.4.6}{}}
\newlabel{eje:4-diofantica@cref}{{[ejercicio][6][4]4.6}{[1][20][]20}}
\newlabel{eje:4-mod4}{{4.7}{20}{}{ejercicio.4.7}{}}
\newlabel{eje:4-mod4@cref}{{[ejercicio][7][4]4.7}{[1][20][]20}}
\newlabel{eje:4-24}{{4.8}{20}{}{ejercicio.4.8}{}}
\newlabel{eje:4-24@cref}{{[ejercicio][8][4]4.8}{[1][20][]20}}
\newlabel{eje:4-maximo}{{4.9}{20}{}{ejercicio.4.9}{}}
\newlabel{eje:4-maximo@cref}{{[ejercicio][9][4]4.9}{[1][20][]20}}
\newlabel{eje:4-restos-cuadrados}{{4.10}{20}{}{ejercicio.4.10}{}}
\newlabel{eje:4-restos-cuadrados@cref}{{[ejercicio][10][4]4.10}{[1][20][]20}}
\newlabel{eje:4-triangular-inversa}{{4.11}{20}{}{ejercicio.4.11}{}}
\newlabel{eje:4-triangular-inversa@cref}{{[ejercicio][11][4]4.11}{[1][20][]20}}
\newlabel{eje:4-producto-par}{{4.12}{20}{}{ejercicio.4.12}{}}
\newlabel{eje:4-producto-par@cref}{{[ejercicio][12][4]4.12}{[1][20][]20}}
\newlabel{eje:4-cubos}{{4.13}{20}{}{ejercicio.4.13}{}}
\newlabel{eje:4-cubos@cref}{{[ejercicio][13][4]4.13}{[1][20][]20}}
\newlabel{eje:4-euclides-n}{{4.14}{20}{}{ejercicio.4.14}{}}
\newlabel{eje:4-euclides-n@cref}{{[ejercicio][14][4]4.14}{[1][20][]20}}
\newlabel{eje:4-4k3}{{4.15}{20}{}{ejercicio.4.15}{}}
\newlabel{eje:4-4k3@cref}{{[ejercicio][15][4]4.15}{[1][20][]20}}
\newlabel{eje:4-raiz-primo}{{4.16}{20}{}{ejercicio.4.16}{}}
\newlabel{eje:4-raiz-primo@cref}{{[ejercicio][16][4]4.16}{[1][20][]20}}
\newlabel{ej:gauss}{{5.2}{22}{Suma de los primeros naturales}{ejemplo.2}{}}
\newlabel{ej:gauss@cref}{{[ejemplo][2][]5.2}{[1][22][]22}}
\newlabel{ej:geometrica}{{5.3}{22}{Serie geométrica finita}{ejemplo.3}{}}
\newlabel{ej:geometrica@cref}{{[ejemplo][3][]5.3}{[1][22][]22}}
\newlabel{ej:desig-induccion}{{5.4}{23}{Una desigualdad}{ejemplo.4}{}}
\newlabel{ej:desig-induccion@cref}{{[ejemplo][4][]5.4}{[1][22][]23}}
\newlabel{ej:div-induccion}{{5.5}{23}{Divisibilidad}{ejemplo.5}{}}
\newlabel{ej:div-induccion@cref}{{[ejemplo][5][]5.5}{[1][23][]23}}
\newlabel{ej:sellos}{{5.8}{23}{Varios casos base}{ejemplo.8}{}}
\newlabel{ej:sellos@cref}{{[ejemplo][8][]5.8}{[1][23][]23}}
\newlabel{obs:5-errores}{{5.12}{24}{Errores frecuentes}{observacion.12}{}}
\newlabel{obs:5-errores@cref}{{[observacion][12][]5.12}{[1][24][]24}}
\newlabel{eje:5-plantilla}{{5.1}{25}{}{ejercicio.5.1}{}}
\newlabel{eje:5-plantilla@cref}{{[ejercicio][1][5]5.1}{[1][25][]25}}
\newlabel{eje:5-cuadrados}{{5.2}{25}{}{ejercicio.5.2}{}}
\newlabel{eje:5-cuadrados@cref}{{[ejercicio][2][5]5.2}{[1][25][]25}}
\newlabel{eje:5-cubos}{{5.3}{25}{}{ejercicio.5.3}{}}
\newlabel{eje:5-cubos@cref}{{[ejercicio][3][5]5.3}{[1][25][]25}}
\newlabel{eje:5-telescopica}{{5.4}{25}{}{ejercicio.5.4}{}}
\newlabel{eje:5-telescopica@cref}{{[ejercicio][4][5]5.4}{[1][25][]25}}
\newlabel{eje:5-potencia}{{5.5}{25}{}{ejercicio.5.5}{}}
\newlabel{eje:5-potencia@cref}{{[ejercicio][5][5]5.5}{[1][25][]25}}
\newlabel{eje:5-div9}{{5.6}{25}{}{ejercicio.5.6}{}}
\newlabel{eje:5-div9@cref}{{[ejercicio][6][5]5.6}{[1][25][]25}}
\newlabel{eje:5-bernoulli}{{5.7}{25}{}{ejercicio.5.7}{}}
\newlabel{eje:5-bernoulli@cref}{{[ejercicio][7][5]5.7}{[1][25][]25}}
\newlabel{eje:5-sellos}{{5.8}{25}{}{ejercicio.5.8}{}}
\newlabel{eje:5-sellos@cref}{{[ejercicio][8][5]5.8}{[1][25][]25}}
\newlabel{eje:5-fib-suma}{{5.9}{25}{}{ejercicio.5.9}{}}
\newlabel{eje:5-fib-suma@cref}{{[ejercicio][9][5]5.9}{[1][25][]25}}
\newlabel{eje:5-fib-cota}{{5.10}{25}{}{ejercicio.5.10}{}}
\newlabel{eje:5-fib-cota@cref}{{[ejercicio][10][5]5.10}{[1][25][]25}}
\newlabel{eje:5-binario}{{5.11}{26}{}{ejercicio.5.11}{}}
\newlabel{eje:5-binario@cref}{{[ejercicio][11][5]5.11}{[1][26][]26}}
\newlabel{eje:5-bo-geometrica}{{5.12}{26}{}{ejercicio.5.12}{}}
\newlabel{eje:5-bo-geometrica@cref}{{[ejercicio][12][5]5.12}{[1][26][]26}}
\newlabel{eje:5-caballos}{{5.13}{26}{}{ejercicio.5.13}{}}
\newlabel{eje:5-caballos@cref}{{[ejercicio][13][5]5.13}{[1][26][]26}}
\newlabel{eje:5-factorial}{{5.14}{26}{}{ejercicio.5.14}{}}
\newlabel{eje:5-factorial@cref}{{[ejercicio][14][5]5.14}{[1][26][]26}}
\newlabel{eje:5-sumas}{{5.15}{26}{}{ejercicio.5.15}{}}
\newlabel{eje:5-sumas@cref}{{[ejercicio][15][5]5.15}{[1][26][]26}}
\newlabel{eje:5-log}{{5.16}{26}{}{ejercicio.5.16}{}}
\newlabel{eje:5-log@cref}{{[ejercicio][16][5]5.16}{[1][26][]26}}
\newlabel{ej:1729}{{6.1}{27}{Existencia constructiva}{ejemplo.1}{}}
\newlabel{ej:1729@cref}{{[ejemplo][1][]6.1}{[1][27][]27}}
\newlabel{ej:unico-real}{{6.4}{28}{}{ejemplo.4}{}}
\newlabel{ej:unico-real@cref}{{[ejemplo][4][]6.4}{[1][28][]28}}
\newlabel{ej:euler}{{6.5}{28}{El polinomio de Euler}{ejemplo.5}{}}
\newlabel{ej:euler@cref}{{[ejemplo][5][]6.5}{[1][28][]28}}
\newlabel{ej:divide-producto}{{6.6}{28}{}{ejemplo.6}{}}
\newlabel{ej:divide-producto@cref}{{[ejemplo][6][]6.6}{[1][28][]28}}
\newlabel{ej:redaccion}{{6.7}{29}{Borrador y redacción final}{ejemplo.7}{}}
\newlabel{ej:redaccion@cref}{{[ejemplo][7][]6.7}{[1][29][]29}}
\newlabel{obs:6-errores}{{6.8}{29}{Errores frecuentes}{observacion.8}{}}
\newlabel{obs:6-errores@cref}{{[observacion][8][]6.8}{[1][29][]29}}
\newlabel{eje:6-clasificar}{{6.1}{29}{}{ejercicio.6.1}{}}
\newlabel{eje:6-clasificar@cref}{{[ejercicio][1][6]6.1}{[1][29][]29}}
\newlabel{eje:6-contraejemplos}{{6.2}{29}{}{ejercicio.6.2}{}}
\newlabel{eje:6-contraejemplos@cref}{{[ejercicio][2][6]6.2}{[1][29][]29}}
\newlabel{eje:6-demostrar-refutar}{{6.3}{30}{}{ejercicio.6.3}{}}
\newlabel{eje:6-demostrar-refutar@cref}{{[ejercicio][3][6]6.3}{[1][30][]30}}
\newlabel{eje:6-testigos}{{6.4}{30}{}{ejercicio.6.4}{}}
\newlabel{eje:6-testigos@cref}{{[ejercicio][4][6]6.4}{[1][30][]30}}
\newlabel{eje:6-irracional-entre}{{6.5}{30}{}{ejercicio.6.5}{}}
\newlabel{eje:6-irracional-entre@cref}{{[ejercicio][5][6]6.5}{[1][30][]30}}
\newlabel{eje:6-paridad-unica}{{6.6}{30}{}{ejercicio.6.6}{}}
\newlabel{eje:6-paridad-unica@cref}{{[ejercicio][6][6]6.6}{[1][30][]30}}
\newlabel{eje:6-terna}{{6.7}{30}{}{ejercicio.6.7}{}}
\newlabel{eje:6-terna@cref}{{[ejercicio][7][6]6.7}{[1][30][]30}}
\newlabel{eje:6-mersenne}{{6.8}{30}{}{ejercicio.6.8}{}}
\newlabel{eje:6-mersenne@cref}{{[ejercicio][8][6]6.8}{[1][30][]30}}
\newlabel{eje:6-promedio}{{6.9}{30}{}{ejercicio.6.9}{}}
\newlabel{eje:6-promedio@cref}{{[ejercicio][9][6]6.9}{[1][30][]30}}
\newlabel{eje:6-critica-1}{{6.10}{30}{}{ejercicio.6.10}{}}
\newlabel{eje:6-critica-1@cref}{{[ejercicio][10][6]6.10}{[1][30][]30}}
\newlabel{eje:6-critica-2}{{6.11}{30}{}{ejercicio.6.11}{}}
\newlabel{eje:6-critica-2@cref}{{[ejercicio][11][6]6.11}{[1][30][]30}}
\newlabel{eje:6-uno-dos}{{6.12}{30}{}{ejercicio.6.12}{}}
\newlabel{eje:6-uno-dos@cref}{{[ejercicio][12][6]6.12}{[1][30][]30}}
\newlabel{eje:6-redactar}{{6.13}{30}{}{ejercicio.6.13}{}}
\newlabel{eje:6-redactar@cref}{{[ejercicio][13][6]6.13}{[1][30][]30}}
\newlabel{eje:6-unica-raiz}{{6.14}{31}{}{ejercicio.6.14}{}}
\newlabel{eje:6-unica-raiz@cref}{{[ejercicio][14][6]6.14}{[1][31][]31}}
\newlabel{ap:soluciones}{{A}{32}{Soluciones de los ejercicios}{appendix.A}{}}
\newlabel{ap:soluciones@cref}{{[appendix][1][2147483647]A}{[1][32][]32}}
\newlabel{sol:eje:1-proposiciones}{{A.1}{32}{Tema 1}{section*.3}{}}
\newlabel{sol:eje:1-proposiciones@cref}{{[subappendix][1][2147483647,1]A.1}{[1][32][]32}}
\newlabel{sol:eje:1-traducir}{{A.1}{32}{Tema 1}{section*.4}{}}
\newlabel{sol:eje:1-traducir@cref}{{[subappendix][1][2147483647,1]A.1}{[1][32][]32}}
\newlabel{sol:eje:1-reciproco}{{A.1}{33}{Tema 1}{section*.5}{}}
\newlabel{sol:eje:1-reciproco@cref}{{[subappendix][1][2147483647,1]A.1}{[1][33][]33}}
\newlabel{sol:eje:1-necesaria}{{A.1}{33}{Tema 1}{section*.6}{}}
\newlabel{sol:eje:1-necesaria@cref}{{[subappendix][1][2147483647,1]A.1}{[1][33][]33}}
\newlabel{sol:eje:1-silogismo}{{A.1}{33}{Tema 1}{section*.7}{}}
\newlabel{sol:eje:1-silogismo@cref}{{[subappendix][1][2147483647,1]A.1}{[1][33][]33}}
\newlabel{sol:eje:1-xor}{{A.1}{34}{Tema 1}{section*.8}{}}
\newlabel{sol:eje:1-xor@cref}{{[subappendix][1][2147483647,1]A.1}{[1][34][]34}}
\newlabel{sol:eje:1-equivalencias}{{A.1}{34}{Tema 1}{section*.9}{}}
\newlabel{sol:eje:1-equivalencias@cref}{{[subappendix][1][2147483647,1]A.1}{[1][34][]34}}
\newlabel{sol:eje:1-algebra}{{A.1}{35}{Tema 1}{section*.10}{}}
\newlabel{sol:eje:1-algebra@cref}{{[subappendix][1][2147483647,1]A.1}{[1][35][]35}}
\newlabel{sol:eje:1-validez}{{A.1}{35}{Tema 1}{section*.11}{}}
\newlabel{sol:eje:1-validez@cref}{{[subappendix][1][2147483647,1]A.1}{[1][35][]35}}
\newlabel{sol:eje:1-contar}{{A.1}{36}{Tema 1}{section*.12}{}}
\newlabel{sol:eje:1-contar@cref}{{[subappendix][1][2147483647,1]A.1}{[1][35][]36}}
\newlabel{sol:eje:1-leyes}{{A.1}{36}{Tema 1}{section*.13}{}}
\newlabel{sol:eje:1-leyes@cref}{{[subappendix][1][2147483647,1]A.1}{[1][36][]36}}
\newlabel{sol:eje:1-demorgan-n}{{A.1}{36}{Tema 1}{section*.14}{}}
\newlabel{sol:eje:1-demorgan-n@cref}{{[subappendix][1][2147483647,1]A.1}{[1][36][]36}}
\newlabel{sol:eje:1-completitud}{{A.1}{37}{Tema 1}{section*.15}{}}
\newlabel{sol:eje:1-completitud@cref}{{[subappendix][1][2147483647,1]A.1}{[1][37][]37}}
\newlabel{sol:eje:1-bicond-asoc}{{A.1}{37}{Tema 1}{section*.16}{}}
\newlabel{sol:eje:1-bicond-asoc@cref}{{[subappendix][1][2147483647,1]A.1}{[1][37][]37}}
\newlabel{sol:eje:2-valores}{{A.2}{38}{Tema 2}{section*.17}{}}
\newlabel{sol:eje:2-valores@cref}{{[subappendix][2][2147483647,1]A.2}{[1][38][]38}}
\newlabel{sol:eje:2-simbolos}{{A.2}{38}{Tema 2}{section*.18}{}}
\newlabel{sol:eje:2-simbolos@cref}{{[subappendix][2][2147483647,1]A.2}{[1][38][]38}}
\newlabel{sol:eje:2-lenguaje}{{A.2}{39}{Tema 2}{section*.19}{}}
\newlabel{sol:eje:2-lenguaje@cref}{{[subappendix][2][2147483647,1]A.2}{[1][39][]39}}
\newlabel{sol:eje:2-negar}{{A.2}{39}{Tema 2}{section*.20}{}}
\newlabel{sol:eje:2-negar@cref}{{[subappendix][2][2147483647,1]A.2}{[1][39][]39}}
\newlabel{sol:eje:2-continuidad}{{A.2}{39}{Tema 2}{section*.21}{}}
\newlabel{sol:eje:2-continuidad@cref}{{[subappendix][2][2147483647,1]A.2}{[1][39][]39}}
\newlabel{sol:eje:2-orden}{{A.2}{40}{Tema 2}{section*.22}{}}
\newlabel{sol:eje:2-orden@cref}{{[subappendix][2][2147483647,1]A.2}{[1][40][]40}}
\newlabel{sol:eje:2-distribuir}{{A.2}{40}{Tema 2}{section*.23}{}}
\newlabel{sol:eje:2-distribuir@cref}{{[subappendix][2][2147483647,1]A.2}{[1][40][]40}}
\newlabel{sol:eje:2-distribuir-dem}{{A.2}{40}{Tema 2}{section*.24}{}}
\newlabel{sol:eje:2-distribuir-dem@cref}{{[subappendix][2][2147483647,1]A.2}{[1][40][]40}}
\newlabel{sol:eje:2-vacio}{{A.2}{41}{Tema 2}{section*.25}{}}
\newlabel{sol:eje:2-vacio@cref}{{[subappendix][2][2147483647,1]A.2}{[1][41][]41}}
\newlabel{sol:eje:2-unicidad}{{A.2}{41}{Tema 2}{section*.26}{}}
\newlabel{sol:eje:2-unicidad@cref}{{[subappendix][2][2147483647,1]A.2}{[1][41][]41}}
\newlabel{sol:eje:2-acotada}{{A.2}{42}{Tema 2}{section*.27}{}}
\newlabel{sol:eje:2-acotada@cref}{{[subappendix][2][2147483647,1]A.2}{[1][41][]42}}
\newlabel{sol:eje:2-arquimedes}{{A.2}{42}{Tema 2}{section*.28}{}}
\newlabel{sol:eje:2-arquimedes@cref}{{[subappendix][2][2147483647,1]A.2}{[1][42][]42}}
\newlabel{sol:eje:2-critica}{{A.2}{42}{Tema 2}{section*.29}{}}
\newlabel{sol:eje:2-critica@cref}{{[subappendix][2][2147483647,1]A.2}{[1][42][]42}}
\newlabel{sol:eje:2-finito}{{A.2}{43}{Tema 2}{section*.30}{}}
\newlabel{sol:eje:2-finito@cref}{{[subappendix][2][2147483647,1]A.2}{[1][42][]43}}
\newlabel{sol:eje:3-hipotesis}{{A.3}{43}{Tema 3}{section*.31}{}}
\newlabel{sol:eje:3-hipotesis@cref}{{[subappendix][3][2147483647,1]A.3}{[1][43][]43}}
\newlabel{sol:eje:3-esquema}{{A.3}{44}{Tema 3}{section*.32}{}}
\newlabel{sol:eje:3-esquema@cref}{{[subappendix][3][2147483647,1]A.3}{[1][43][]44}}
\newlabel{sol:eje:3-cubo}{{A.3}{44}{Tema 3}{section*.33}{}}
\newlabel{sol:eje:3-cubo@cref}{{[subappendix][3][2147483647,1]A.3}{[1][44][]44}}
\newlabel{sol:eje:3-polinomio}{{A.3}{44}{Tema 3}{section*.34}{}}
\newlabel{sol:eje:3-polinomio@cref}{{[subappendix][3][2147483647,1]A.3}{[1][44][]44}}
\newlabel{sol:eje:3-cuadrados}{{A.3}{44}{Tema 3}{section*.35}{}}
\newlabel{sol:eje:3-cuadrados@cref}{{[subappendix][3][2147483647,1]A.3}{[1][44][]44}}
\newlabel{sol:eje:3-producto}{{A.3}{45}{Tema 3}{section*.36}{}}
\newlabel{sol:eje:3-producto@cref}{{[subappendix][3][2147483647,1]A.3}{[1][45][]45}}
\newlabel{sol:eje:3-desigualdad}{{A.3}{45}{Tema 3}{section*.37}{}}
\newlabel{sol:eje:3-desigualdad@cref}{{[subappendix][3][2147483647,1]A.3}{[1][45][]45}}
\newlabel{sol:eje:3-diferencia}{{A.3}{45}{Tema 3}{section*.38}{}}
\newlabel{sol:eje:3-diferencia@cref}{{[subappendix][3][2147483647,1]A.3}{[1][45][]45}}
\newlabel{sol:eje:3-contra-cubo}{{A.3}{46}{Tema 3}{section*.39}{}}
\newlabel{sol:eje:3-contra-cubo@cref}{{[subappendix][3][2147483647,1]A.3}{[1][46][]46}}
\newlabel{sol:eje:3-pitagoras}{{A.3}{46}{Tema 3}{section*.40}{}}
\newlabel{sol:eje:3-pitagoras@cref}{{[subappendix][3][2147483647,1]A.3}{[1][46][]46}}
\newlabel{sol:eje:3-equivalencia}{{A.3}{46}{Tema 3}{section*.41}{}}
\newlabel{sol:eje:3-equivalencia@cref}{{[subappendix][3][2147483647,1]A.3}{[1][46][]46}}
\newlabel{sol:eje:3-real}{{A.3}{47}{Tema 3}{section*.42}{}}
\newlabel{sol:eje:3-real@cref}{{[subappendix][3][2147483647,1]A.3}{[1][47][]47}}
\newlabel{sol:eje:3-bicond}{{A.3}{47}{Tema 3}{section*.43}{}}
\newlabel{sol:eje:3-bicond@cref}{{[subappendix][3][2147483647,1]A.3}{[1][47][]47}}
\newlabel{sol:eje:3-restos}{{A.3}{47}{Tema 3}{section*.44}{}}
\newlabel{sol:eje:3-restos@cref}{{[subappendix][3][2147483647,1]A.3}{[1][47][]47}}
\newlabel{sol:eje:3-division}{{A.3}{48}{Tema 3}{section*.45}{}}
\newlabel{sol:eje:3-division@cref}{{[subappendix][3][2147483647,1]A.3}{[1][48][]48}}
\newlabel{sol:eje:4-suponer}{{A.4}{48}{Tema 4}{section*.46}{}}
\newlabel{sol:eje:4-suponer@cref}{{[subappendix][4][2147483647,1]A.4}{[1][48][]48}}
\newlabel{sol:eje:4-raiz3}{{A.4}{49}{Tema 4}{section*.47}{}}
\newlabel{sol:eje:4-raiz3@cref}{{[subappendix][4][2147483647,1]A.4}{[1][49][]49}}
\newlabel{sol:eje:4-raiz-cubica}{{A.4}{49}{Tema 4}{section*.48}{}}
\newlabel{sol:eje:4-raiz-cubica@cref}{{[subappendix][4][2147483647,1]A.4}{[1][49][]49}}
\newlabel{sol:eje:4-log}{{A.4}{49}{Tema 4}{section*.49}{}}
\newlabel{sol:eje:4-log@cref}{{[subappendix][4][2147483647,1]A.4}{[1][49][]49}}
\newlabel{sol:eje:4-racional-irracional}{{A.4}{50}{Tema 4}{section*.50}{}}
\newlabel{sol:eje:4-racional-irracional@cref}{{[subappendix][4][2147483647,1]A.4}{[1][50][]50}}
\newlabel{sol:eje:4-diofantica}{{A.4}{50}{Tema 4}{section*.51}{}}
\newlabel{sol:eje:4-diofantica@cref}{{[subappendix][4][2147483647,1]A.4}{[1][50][]50}}
\newlabel{sol:eje:4-mod4}{{A.4}{51}{Tema 4}{section*.52}{}}
\newlabel{sol:eje:4-mod4@cref}{{[subappendix][4][2147483647,1]A.4}{[1][50][]51}}
\newlabel{sol:eje:4-24}{{A.4}{51}{Tema 4}{section*.53}{}}
\newlabel{sol:eje:4-24@cref}{{[subappendix][4][2147483647,1]A.4}{[1][51][]51}}
\newlabel{sol:eje:4-maximo}{{A.4}{51}{Tema 4}{section*.54}{}}
\newlabel{sol:eje:4-maximo@cref}{{[subappendix][4][2147483647,1]A.4}{[1][51][]51}}
\newlabel{sol:eje:4-restos-cuadrados}{{A.4}{52}{Tema 4}{section*.55}{}}
\newlabel{sol:eje:4-restos-cuadrados@cref}{{[subappendix][4][2147483647,1]A.4}{[1][52][]52}}
\newlabel{sol:eje:4-triangular-inversa}{{A.4}{52}{Tema 4}{section*.56}{}}
\newlabel{sol:eje:4-triangular-inversa@cref}{{[subappendix][4][2147483647,1]A.4}{[1][52][]52}}
\newlabel{sol:eje:4-producto-par}{{A.4}{52}{Tema 4}{section*.57}{}}
\newlabel{sol:eje:4-producto-par@cref}{{[subappendix][4][2147483647,1]A.4}{[1][52][]52}}
\newlabel{sol:eje:4-cubos}{{A.4}{53}{Tema 4}{section*.58}{}}
\newlabel{sol:eje:4-cubos@cref}{{[subappendix][4][2147483647,1]A.4}{[1][53][]53}}
\newlabel{sol:eje:4-euclides-n}{{A.4}{53}{Tema 4}{section*.59}{}}
\newlabel{sol:eje:4-euclides-n@cref}{{[subappendix][4][2147483647,1]A.4}{[1][53][]53}}
\newlabel{sol:eje:4-4k3}{{A.4}{54}{Tema 4}{section*.60}{}}
\newlabel{sol:eje:4-4k3@cref}{{[subappendix][4][2147483647,1]A.4}{[1][54][]54}}
\newlabel{sol:eje:4-raiz-primo}{{A.4}{54}{Tema 4}{section*.61}{}}
\newlabel{sol:eje:4-raiz-primo@cref}{{[subappendix][4][2147483647,1]A.4}{[1][54][]54}}
\newlabel{sol:eje:5-plantilla}{{A.5}{55}{Tema 5}{section*.62}{}}
\newlabel{sol:eje:5-plantilla@cref}{{[subappendix][5][2147483647,1]A.5}{[1][55][]55}}
\newlabel{sol:eje:5-cuadrados}{{A.5}{55}{Tema 5}{section*.63}{}}
\newlabel{sol:eje:5-cuadrados@cref}{{[subappendix][5][2147483647,1]A.5}{[1][55][]55}}
\newlabel{sol:eje:5-cubos}{{A.5}{55}{Tema 5}{section*.64}{}}
\newlabel{sol:eje:5-cubos@cref}{{[subappendix][5][2147483647,1]A.5}{[1][55][]55}}
\newlabel{sol:eje:5-telescopica}{{A.5}{56}{Tema 5}{section*.65}{}}
\newlabel{sol:eje:5-telescopica@cref}{{[subappendix][5][2147483647,1]A.5}{[1][56][]56}}
\newlabel{sol:eje:5-potencia}{{A.5}{56}{Tema 5}{section*.66}{}}
\newlabel{sol:eje:5-potencia@cref}{{[subappendix][5][2147483647,1]A.5}{[1][56][]56}}
\newlabel{sol:eje:5-div9}{{A.5}{57}{Tema 5}{section*.67}{}}
\newlabel{sol:eje:5-div9@cref}{{[subappendix][5][2147483647,1]A.5}{[1][57][]57}}
\newlabel{sol:eje:5-bernoulli}{{A.5}{57}{Tema 5}{section*.68}{}}
\newlabel{sol:eje:5-bernoulli@cref}{{[subappendix][5][2147483647,1]A.5}{[1][57][]57}}
\newlabel{sol:eje:5-sellos}{{A.5}{57}{Tema 5}{section*.69}{}}
\newlabel{sol:eje:5-sellos@cref}{{[subappendix][5][2147483647,1]A.5}{[1][57][]57}}
\newlabel{sol:eje:5-fib-suma}{{A.5}{58}{Tema 5}{section*.70}{}}
\newlabel{sol:eje:5-fib-suma@cref}{{[subappendix][5][2147483647,1]A.5}{[1][58][]58}}
\newlabel{sol:eje:5-fib-cota}{{A.5}{58}{Tema 5}{section*.71}{}}
\newlabel{sol:eje:5-fib-cota@cref}{{[subappendix][5][2147483647,1]A.5}{[1][58][]58}}
\newlabel{sol:eje:5-binario}{{A.5}{58}{Tema 5}{section*.72}{}}
\newlabel{sol:eje:5-binario@cref}{{[subappendix][5][2147483647,1]A.5}{[1][58][]58}}
\newlabel{sol:eje:5-bo-geometrica}{{A.5}{59}{Tema 5}{section*.73}{}}
\newlabel{sol:eje:5-bo-geometrica@cref}{{[subappendix][5][2147483647,1]A.5}{[1][59][]59}}
\newlabel{sol:eje:5-caballos}{{A.5}{59}{Tema 5}{section*.74}{}}
\newlabel{sol:eje:5-caballos@cref}{{[subappendix][5][2147483647,1]A.5}{[1][59][]59}}
\newlabel{sol:eje:5-factorial}{{A.5}{60}{Tema 5}{section*.75}{}}
\newlabel{sol:eje:5-factorial@cref}{{[subappendix][5][2147483647,1]A.5}{[1][60][]60}}
\newlabel{sol:eje:5-sumas}{{A.5}{60}{Tema 5}{section*.76}{}}
\newlabel{sol:eje:5-sumas@cref}{{[subappendix][5][2147483647,1]A.5}{[1][60][]60}}
\newlabel{sol:eje:5-log}{{A.5}{60}{Tema 5}{section*.77}{}}
\newlabel{sol:eje:5-log@cref}{{[subappendix][5][2147483647,1]A.5}{[1][60][]60}}
\newlabel{sol:eje:6-clasificar}{{A.6}{61}{Tema 6}{section*.78}{}}
\newlabel{sol:eje:6-clasificar@cref}{{[subappendix][6][2147483647,1]A.6}{[1][61][]61}}
\newlabel{sol:eje:6-contraejemplos}{{A.6}{61}{Tema 6}{section*.79}{}}
\newlabel{sol:eje:6-contraejemplos@cref}{{[subappendix][6][2147483647,1]A.6}{[1][61][]61}}
\newlabel{sol:eje:6-demostrar-refutar}{{A.6}{62}{Tema 6}{section*.80}{}}
\newlabel{sol:eje:6-demostrar-refutar@cref}{{[subappendix][6][2147483647,1]A.6}{[1][62][]62}}
\newlabel{sol:eje:6-testigos}{{A.6}{62}{Tema 6}{section*.81}{}}
\newlabel{sol:eje:6-testigos@cref}{{[subappendix][6][2147483647,1]A.6}{[1][62][]62}}
\newlabel{sol:eje:6-irracional-entre}{{A.6}{62}{Tema 6}{section*.82}{}}
\newlabel{sol:eje:6-irracional-entre@cref}{{[subappendix][6][2147483647,1]A.6}{[1][62][]62}}
\newlabel{sol:eje:6-paridad-unica}{{A.6}{63}{Tema 6}{section*.83}{}}
\newlabel{sol:eje:6-paridad-unica@cref}{{[subappendix][6][2147483647,1]A.6}{[1][63][]63}}
\newlabel{sol:eje:6-terna}{{A.6}{63}{Tema 6}{section*.84}{}}
\newlabel{sol:eje:6-terna@cref}{{[subappendix][6][2147483647,1]A.6}{[1][63][]63}}
\newlabel{sol:eje:6-mersenne}{{A.6}{64}{Tema 6}{section*.85}{}}
\newlabel{sol:eje:6-mersenne@cref}{{[subappendix][6][2147483647,1]A.6}{[1][63][]64}}
\newlabel{sol:eje:6-promedio}{{A.6}{64}{Tema 6}{section*.86}{}}
\newlabel{sol:eje:6-promedio@cref}{{[subappendix][6][2147483647,1]A.6}{[1][64][]64}}
\newlabel{sol:eje:6-critica-1}{{A.6}{64}{Tema 6}{section*.87}{}}
\newlabel{sol:eje:6-critica-1@cref}{{[subappendix][6][2147483647,1]A.6}{[1][64][]64}}
\newlabel{sol:eje:6-critica-2}{{A.6}{65}{Tema 6}{section*.88}{}}
\newlabel{sol:eje:6-critica-2@cref}{{[subappendix][6][2147483647,1]A.6}{[1][65][]65}}
\newlabel{sol:eje:6-uno-dos}{{A.6}{65}{Tema 6}{section*.89}{}}
\newlabel{sol:eje:6-uno-dos@cref}{{[subappendix][6][2147483647,1]A.6}{[1][65][]65}}
\newlabel{sol:eje:6-redactar}{{A.6}{65}{Tema 6}{section*.90}{}}
\newlabel{sol:eje:6-redactar@cref}{{[subappendix][6][2147483647,1]A.6}{[1][65][]65}}
\newlabel{sol:eje:6-unica-raiz}{{A.6}{66}{Tema 6}{section*.91}{}}
\newlabel{sol:eje:6-unica-raiz@cref}{{[subappendix][6][2147483647,1]A.6}{[1][66][]66}}
\makeatother

\begin{document}
\frontmatter
\begin{titlepage}
\centering
\vspace*{2.5cm}
{\large\scshape Anatomía de un LLM}\par
\vspace{0.4cm}
{\normalsize Fundamentos matemáticos de los modelos de lenguaje}\par
\vspace{2.2cm}
\rule{\textwidth}{0.8pt}\par\vspace{0.5cm}
{\small\scshape Fase 0 \textperiodcentered\ Preliminares matemáticos}\par\vspace{0.3cm}
{\Huge\bfseries Módulo 0.1\par}\vspace{0.3cm}
{\LARGE Lógica y técnicas de demostración\par}
\vspace{0.5cm}\rule{\textwidth}{0.4pt}\par
\vspace{1.6cm}
{\large Resumen del módulo}\par
\vfill
{\small Edición de octubre de 2026}\par
\end{titlepage}

\chapter*{Sobre este resumen}
\addcontentsline{toc}{chapter}{Sobre este resumen}
\markboth{Sobre este resumen}{}

Este documento reúne lo indispensable de los apuntes del Módulo 0.1: el texto expositivo, todas las definiciones y todos los resultados (teoremas, proposiciones, lemas y corolarios), las tablas y las observaciones de contenido. Se han suprimido las demostraciones, los ejemplos, las lecturas, las figuras, los ejercicios y sus soluciones. La numeración de temas, secciones, enunciados y ecuaciones es la misma que en los apuntes, de modo que cualquier referencia puede buscarse allí; las que remiten a elementos suprimidos (un ejemplo, una figura, una ecuación de una demostración) conservan su número en los apuntes.

\section*{Notación}
\begin{center}
\begin{tabular}{@{}ll@{}}
\toprule
Símbolo & Significado \\
\midrule
$\neg P$, $P\wedge Q$, $P\vee Q$ & negación, conjunción y disyunción (inclusiva) \\
$P\Rightarrow Q$, $P\Leftrightarrow Q$ & condicional («si $P$, entonces $Q$») y bicondicional \\
$\V$, $\F$ & valores de verdad «verdadero» y «falso» \\
$\Phi\eqlog\Psi$ & equivalencia lógica de dos fórmulas \\
$\forall$, $\exists$, $\exists!$ & «para todo», «existe», «existe un único» \\
$\N=\{1,2,3,\dots\}$, $\N_0=\N\cup\{0\}$ & números naturales, sin y con el cero \\
$\Z$, $\Q$, $\R$ & enteros, racionales y reales \\
$a\divides b$, $a\nodivide b$ & $a$ divide a $b$; $a$ no divide a $b$ \\
$a\equiv b\pmod n$ & $a$ y $b$ son congruentes módulo $n$ \\
$\floor{x}$, $\ceil{x}$ & parte entera por defecto y por exceso \\
$\blacksquare$ & fin de una demostración \\
\bottomrule
\end{tabular}
\end{center}
El símbolo $\equiv$ se usa con dos sentidos que el contexto distingue: entre fórmulas lógicas indica equivalencia lógica, y en la notación $a\equiv b\pmod n$ indica congruencia.


\tableofcontents

\mainmatter
% =====================================================================
\chapter{Proposiciones, conectivas y tablas de verdad}
\label{tema:1}
% =====================================================================

Una demostración es una cadena de afirmaciones enlazadas por reglas de razonamiento. Antes de estudiar las técnicas de demostración conviene fijar qué es una afirmación, cómo se combinan varias mediante conectivas y cuándo dos combinaciones dicen lo mismo. Este tema desarrolla la lógica proposicional con el enfoque del capítulo~2 de~\cite{hammack} (§2.1–2.6 y §2.11) y del capítulo~1 de~\cite{velleman}; la parte algebraica sigue la §3.4 de~\cite{mcs}.

\section{Proposiciones y conectivas}

\begin{definicion}[Proposición]\label{def:proposicion}
Una \emph{proposición} es un enunciado del que tiene sentido afirmar que es verdadero o que es falso, y que es exactamente una de las dos cosas. Su \emph{valor de verdad} es $\V$ (verdadero) o $\F$ (falso).
\end{definicion}

\addtocounter{definicion}{1}%

A partir de proposiciones $P$ y $Q$ se forman proposiciones compuestas mediante \emph{conectivas lógicas}. El valor de verdad de la proposición compuesta depende solo de los de $P$ y $Q$, y se describe con una \emph{tabla de verdad}.

\begin{definicion}[Negación, conjunción y disyunción]\label{def:conectivas}
Sean $P$ y $Q$ proposiciones.
\begin{enumerate}
  \item La \emph{negación} $\neg P$ («no $P$») es verdadera si y solo si $P$ es falsa.
  \item La \emph{conjunción} $P\wedge Q$ («$P$ y $Q$») es verdadera si y solo si $P$ y $Q$ son ambas verdaderas.
  \item La \emph{disyunción} $P\vee Q$ («$P$ o $Q$») es verdadera si y solo si al menos una de las dos es verdadera.
\end{enumerate}
\end{definicion}

\begin{table}[h]
\centering
\caption{Tablas de verdad de las conectivas básicas.}\label{tab:conectivas}
\begin{tabular}{@{}cc|ccccc@{}}
\toprule
$P$ & $Q$ & $\neg P$ & $P\wedge Q$ & $P\vee Q$ & $P\Rightarrow Q$ & $P\Leftrightarrow Q$ \\
\midrule
$\V$ & $\V$ & $\F$ & $\V$ & $\V$ & $\V$ & $\V$ \\
$\V$ & $\F$ & $\F$ & $\F$ & $\V$ & $\F$ & $\F$ \\
$\F$ & $\V$ & $\V$ & $\F$ & $\V$ & $\V$ & $\F$ \\
$\F$ & $\F$ & $\V$ & $\F$ & $\F$ & $\V$ & $\V$ \\
\bottomrule
\end{tabular}
\end{table}

La disyunción matemática es \emph{inclusiva}: $P\vee Q$ es verdadera también cuando lo son las dos. La disyunción \emph{exclusiva}, «o $P$ o $Q$, pero no ambas», se escribe $P\xor Q$ y se estudia en el \cref{eje:1-xor}.

\section{Condicional y bicondicional}

\begin{definicion}[Condicional]\label{def:condicional}
El \emph{condicional} $P\Rightarrow Q$ («si $P$, entonces $Q$») es falso únicamente cuando $P$ es verdadera y $Q$ es falsa. $P$ se llama \emph{hipótesis} o antecedente, y $Q$, \emph{tesis} o consecuente.
\end{definicion}

La tabla del condicional (\cref{tab:conectivas}) se entiende como una promesa: «si apruebas el examen, te invito a cenar» solo se incumple si se aprueba y no hay cena. Si no se aprueba, la promesa no se ha roto, haya cena o no. Por eso un condicional con hipótesis falsa es verdadero; se dice que es \emph{vacuamente} verdadero.

El condicional $P\Rightarrow Q$ admite varias lecturas equivalentes, que conviene reconocer porque aparecen constantemente en los enunciados de los teoremas~\cite[§2.3]{hammack}:
\begin{center}
\begin{tabular}{@{}l@{}}
«si $P$, entonces $Q$»; \quad «$Q$ si $P$»; \quad «$P$ solo si $Q$»; \\
«$P$ es condición suficiente para $Q$»; \quad «$Q$ es condición necesaria para $P$».
\end{tabular}
\end{center}

\begin{definicion}[Bicondicional]\label{def:bicondicional}
El \emph{bicondicional} $P\Leftrightarrow Q$ («$P$ si y solo si $Q$») es verdadero cuando $P$ y $Q$ tienen el mismo valor de verdad. Se lee también «$P$ es condición necesaria y suficiente para $Q$».
\end{definicion}

\begin{definicion}[Recíproco, contrarrecíproco e inverso]\label{def:reciproco}
Dado el condicional $P\Rightarrow Q$, se llaman \emph{recíproco} a $Q\Rightarrow P$, \emph{contrarrecíproco} a $\neg Q\Rightarrow\neg P$ e \emph{inverso} a $\neg P\Rightarrow\neg Q$.
\end{definicion}

\addtocounter{definicion}{1}%

\section{Fórmulas, tautologías y contradicciones}

\begin{definicion}[Fórmula proposicional]
Las \emph{fórmulas proposicionales} sobre las variables $P_1,\dots,P_n$ son las expresiones que se obtienen a partir de las variables aplicando un número finito de veces las conectivas $\neg$, $\wedge$, $\vee$, $\Rightarrow$ y $\Leftrightarrow$, con paréntesis para indicar el orden. Una \emph{asignación} da a cada variable un valor de verdad, y la tabla de verdad de la fórmula lista su valor para cada asignación.
\end{definicion}

\begin{proposicion}\label{prop:filas}
Una fórmula con $n$ variables tiene una tabla de verdad de $2^{n}$ filas.
\end{proposicion}

\begin{definicion}[Tautología y contradicción]
Una fórmula es una \emph{tautología} si es verdadera para todas las asignaciones, y una \emph{contradicción} si es falsa para todas.
\end{definicion}

\addtocounter{definicion}{1}%

El razonamiento del ejemplo anterior evita escribir la tabla completa: para comprobar que un condicional es tautología basta ver que la fila «hipótesis verdadera y tesis falsa» es imposible. Esta forma de pensar es ya una demostración por reducción al absurdo (\cref{tema:4}).

\section{Equivalencia lógica}

\begin{definicion}[Equivalencia lógica]\label{def:equivalencia}
Dos fórmulas $\Phi$ y $\Psi$ en las mismas variables son \emph{lógicamente equivalentes}, y se escribe $\Phi\eqlog\Psi$, si tienen el mismo valor de verdad en todas las asignaciones; equivalentemente, si $\Phi\Leftrightarrow\Psi$ es una tautología.
\end{definicion}

\begin{proposicion}[Leyes lógicas]\label{prop:leyes}
Para proposiciones cualesquiera $P$, $Q$ y $R$:
\begin{enumerate}
  \item doble negación: $\neg\neg P\eqlog P$;
  \item leyes de De Morgan: $\neg(P\wedge Q)\eqlog\neg P\vee\neg Q$ \ y \ $\neg(P\vee Q)\eqlog\neg P\wedge\neg Q$;
  \item distributivas: $P\wedge(Q\vee R)\eqlog(P\wedge Q)\vee(P\wedge R)$ \ y \ $P\vee(Q\wedge R)\eqlog(P\vee Q)\wedge(P\vee R)$;
  \item condicional como disyunción: $P\Rightarrow Q\eqlog\neg P\vee Q$;
  \item contrarrecíproco: $P\Rightarrow Q\eqlog\neg Q\Rightarrow\neg P$;
  \item negación del condicional: $\neg(P\Rightarrow Q)\eqlog P\wedge\neg Q$;
  \item bicondicional: $P\Leftrightarrow Q\eqlog(P\Rightarrow Q)\wedge(Q\Rightarrow P)$.
\end{enumerate}
Además, $\wedge$ y $\vee$ son conmutativas y asociativas.
\end{proposicion}

Las leyes de la \cref{prop:leyes} permiten transformar fórmulas sin construir tablas, igual que las propiedades de la suma y el producto permiten manipular expresiones algebraicas~\cite[§3.4]{mcs}.

\addtocounter{definicion}{1}%

\section{Razonamientos válidos}\label{sec:inferencia}

\begin{definicion}[Razonamiento válido]
Un razonamiento con premisas $\Phi_1,\dots,\Phi_k$ y conclusión $\Psi$ es \emph{válido} si $(\Phi_1\wedge\dots\wedge\Phi_k)\Rightarrow\Psi$ es una tautología, es decir, si en toda asignación que hace verdaderas todas las premisas también es verdadera la conclusión.
\end{definicion}

\begin{proposicion}[Reglas de inferencia]\label{prop:inferencia}
Los siguientes razonamientos son válidos~\cite[§2.11]{hammack}:
\begin{center}
\begin{tabular}{@{}lll@{}}
\toprule
Nombre & Premisas & Conclusión \\
\midrule
\emph{Modus ponens} & $P\Rightarrow Q$, \ $P$ & $Q$ \\
\emph{Modus tollens} & $P\Rightarrow Q$, \ $\neg Q$ & $\neg P$ \\
Silogismo hipotético & $P\Rightarrow Q$, \ $Q\Rightarrow R$ & $P\Rightarrow R$ \\
Silogismo disyuntivo & $P\vee Q$, \ $\neg P$ & $Q$ \\
\bottomrule
\end{tabular}
\end{center}
\end{proposicion}

\addtocounter{definicion}{1}%

\addtocounter{definicion}{1}%


% =====================================================================
\chapter{Cuantificadores y negación}
\label{tema:2}
% =====================================================================

Casi ningún teorema habla de un único objeto: afirman algo de \emph{todos} los enteros, de \emph{alguna} raíz de un polinomio o de que \emph{para cada} $\varepsilon$ \emph{existe} un $\delta$. Este tema introduce los predicados y los cuantificadores, estudia el efecto del orden de los cuantificadores y da reglas sistemáticas para negar enunciados, que serán imprescindibles en las demostraciones por contrarrecíproco y por reducción al absurdo. Se siguen las §2.7–2.10 de~\cite{hammack}, el capítulo~2 de~\cite{velleman} (§2.1 y §2.2), la §3.6 de~\cite{mcs} y la guía de negación del curso CS103~\cite{cs103}.

\section{Predicados}

\begin{definicion}[Predicado]\label{def:predicado}
Sea $U$ un conjunto no vacío, llamado \emph{universo} o dominio. Un \emph{predicado} sobre $U$ es un enunciado $P(x)$ que contiene una variable $x$ y que se convierte en una proposición cada vez que se sustituye $x$ por un elemento concreto de $U$. El \emph{conjunto de verdad} de $P$ es $\{x\in U: P(x)\text{ es verdadera}\}$. Del mismo modo se definen predicados de varias variables, $P(x,y)$, con $(x,y)\in U\times W$.
\end{definicion}

\addtocounter{definicion}{1}%

\section{Cuantificador universal y existencial}

\begin{definicion}[Cuantificadores]\label{def:cuantificadores}
Sea $P(x)$ un predicado sobre $U$.
\begin{enumerate}
  \item $\forall x\in U\ P(x)$ («para todo $x$ de $U$, $P(x)$») es verdadera si $P(a)$ es verdadera para cada $a\in U$.
  \item $\exists x\in U\ P(x)$ («existe $x$ en $U$ tal que $P(x)$») es verdadera si $P(a)$ es verdadera para al menos un $a\in U$.
\end{enumerate}
Cuando el universo está claro se escribe $\forall x\,P(x)$ y $\exists x\,P(x)$. La variable $x$ está \emph{ligada} por el cuantificador: $\forall x\,P(x)$ y $\forall t\,P(t)$ dicen lo mismo.
\end{definicion}

Los cuantificadores se combinan con las conectivas. Las formas más frecuentes son la \emph{universal condicional}, $\forall x\,\bigl(P(x)\Rightarrow Q(x)\bigr)$, que traduce «todo $P$ es $Q$», y la \emph{existencial conjuntiva}, $\exists x\,\bigl(P(x)\wedge Q(x)\bigr)$, que traduce «algún $P$ es $Q$». Se abrevian como $\forall x\in A\ Q(x)$ y $\exists x\in A\ Q(x)$ cuando $P(x)$ es «$x\in A$».

\addtocounter{definicion}{1}%

\begin{observacion}
Si $A=\varnothing$, el enunciado $\forall x\in A\ Q(x)$ es verdadero, sea cual sea $Q$: no hay ningún $x\in A$ que lo incumpla. Es la versión cuantificada del condicional vacuamente verdadero del \cref{tema:1}. En cambio, $\exists x\in\varnothing\ Q(x)$ es siempre falso.
\end{observacion}

\section{Cuantificadores anidados}

Con predicados de varias variables, los cuantificadores se anidan. El orden importa cuando se mezclan cuantificadores de distinto tipo.

\addtocounter{definicion}{1}%

\begin{proposicion}\label{prop:orden}
Para todo predicado $P(x,y)$,
\[
\exists y\,\forall x\,P(x,y)\ \Longrightarrow\ \forall x\,\exists y\,P(x,y),
\]
y la implicación recíproca es falsa en general.
\end{proposicion}

Cuantificadores del mismo tipo sí pueden intercambiarse: $\forall x\,\forall y\,P(x,y)\eqlog\forall y\,\forall x\,P(x,y)$ y $\exists x\,\exists y\,P(x,y)\eqlog\exists y\,\exists x\,P(x,y)$.

\addtocounter{definicion}{1}%

\section{Negación de enunciados cuantificados}

\begin{teorema}[Negación de cuantificadores]\label{teo:negacion}
Para todo predicado $P(x)$ sobre un universo $U$,
\[
\neg\,\forall x\,P(x)\ \eqlog\ \exists x\,\neg P(x),\qquad
\neg\,\exists x\,P(x)\ \eqlog\ \forall x\,\neg P(x).
\]
\end{teorema}

El \cref{teo:negacion} y las leyes de la \cref{prop:leyes} dan un procedimiento mecánico: la negación «atraviesa» los cuantificadores cambiando $\forall$ por $\exists$ y viceversa, y al llegar a la parte proposicional se aplican De Morgan y la negación del condicional, $\neg(P\Rightarrow Q)\eqlog P\wedge\neg Q$. Con universos acotados, el cuantificador conserva su dominio: $\neg\,\forall x\in A\ Q(x)\eqlog\exists x\in A\ \neg Q(x)$.

\addtocounter{definicion}{1}%

\section{Traducción y unicidad}

Las expresiones del lenguaje natural ocultan a menudo los cuantificadores. «Los primos mayores que~2 son impares» es universal; «un cuadrado puede ser par» es existencial; «solo los múltiplos de~6 son múltiplos de~2 y de~3» significa «si $n$ es múltiplo de~2 y de~3, entonces es múltiplo de~6». Antes de demostrar un enunciado conviene escribirlo con cuantificadores explícitos~\cite[§2.9]{hammack}.

\begin{definicion}[Existencia y unicidad]\label{def:unicidad}
$\exists!\,x\,P(x)$ («existe un único $x$ tal que $P(x)$») abrevia
\[
\exists x\,\Bigl(P(x)\wedge\forall y\,\bigl(P(y)\Rightarrow y=x\bigr)\Bigr).
\]
\end{definicion}

Así, $\exists!$ contiene dos afirmaciones: hay al menos un elemento con la propiedad (\emph{existencia}) y dos elementos con la propiedad son iguales (\emph{unicidad}). El \cref{tema:6} explica cómo se demuestra cada parte.

\addtocounter{definicion}{1}%

\addtocounter{definicion}{1}%


% =====================================================================
\chapter{Demostración directa y por contrarrecíproco}
\label{tema:3}
% =====================================================================

La mayoría de los teoremas tienen la forma $\forall x\,(P(x)\Rightarrow Q(x))$. Este tema presenta las dos técnicas básicas para demostrarlos: la demostración directa, que parte de $P(x)$ y llega a $Q(x)$, y la demostración por contrarrecíproco, que parte de $\neg Q(x)$ y llega a $\neg P(x)$. Como banco de pruebas se usa la aritmética de los enteros (paridad, divisibilidad y congruencias), igual que en los capítulos~4 y~5 de~\cite{hammack}, la §3.2 de~\cite{velleman}, las §1.5–1.6 de~\cite{mcs} y el primer \emph{problem set} de CS103~\cite{cs103}.

\section{Definiciones de partida}

Una demostración solo puede usar lo que está definido o demostrado antes. Se admiten como conocidas las propiedades algebraicas habituales de $\Z$ (la suma, la resta y el producto de enteros son enteros; asociatividad, conmutatividad, distributividad) y de $\R$ (incluidas las reglas de las desigualdades), y se fijan las definiciones siguientes~\cite[§4.2]{hammack}.

\begin{definicion}[Paridad]\label{def:paridad}
Un entero $n$ es \emph{par} si existe $k\in\Z$ con $n=2k$, e \emph{impar} si existe $k\in\Z$ con $n=2k+1$. Dos enteros tienen la \emph{misma paridad} si ambos son pares o ambos impares.
\end{definicion}

\begin{definicion}[Divisibilidad]\label{def:divide}
Sean $a,b\in\Z$. Se dice que $a$ \emph{divide} a $b$, y se escribe $a\divides b$, si existe $c\in\Z$ con $b=ac$. Se dice también que $b$ es \emph{múltiplo} de $a$ y que $a$ es \emph{divisor} de $b$. Un entero $p\ge 2$ es \emph{primo} si sus únicos divisores positivos son $1$ y $p$; un entero $n\ge 2$ que no es primo es \emph{compuesto}.
\end{definicion}

\begin{teorema}[División euclídea]\label{teo:division}
Sean $a\in\Z$ y $b\in\N$. Existen enteros únicos $q$ y $r$ con $a=qb+r$ y $0\le r<b$.
\end{teorema}

La existencia se demuestra en el \cref{tema:5} (\cref{prop:division-existencia}) y la unicidad en el \cref{tema:6} (\cref{prop:division-unicidad}). Con $b=2$, el teorema dice que todo entero es par o impar, y no ambas cosas; este hecho se usará libremente.

\section{Demostración directa}

\begin{definicion}[Demostración directa]\label{def:directa}
Para demostrar $\forall x\,(P(x)\Rightarrow Q(x))$ de forma \emph{directa}, se toma un $x$ arbitrario, se supone $P(x)$ y, mediante definiciones, resultados previos y reglas de inferencia, se deduce $Q(x)$.
\end{definicion}

El esquema es siempre el mismo: «Sea $x$ [del universo] tal que $P(x)$. [\dots] Por tanto, $Q(x)$.» Antes de empezar conviene escribir qué significa la tesis según las definiciones, para saber a qué expresión hay que llegar~\cite[§4.3]{hammack}.

\begin{proposicion}\label{prop:impar-cuadrado}
Si $n$ es un entero impar, entonces $n^{2}$ es impar.
\end{proposicion}

Obsérvese que la demostración introduce un nombre nuevo, $m$, para el entero que da la forma $2m+1$: el objetivo era encontrar \emph{algún} entero con esa propiedad, y la definición de impar no exige que sea el mismo $k$ de la hipótesis.

\begin{proposicion}[Propiedades de la divisibilidad]\label{prop:divisibilidad}
Sean $a,b,c,x,y\in\Z$.
\begin{enumerate}
  \item Si $a\divides b$ y $b\divides c$, entonces $a\divides c$.
  \item Si $a\divides b$ y $a\divides c$, entonces $a\divides(bx+cy)$.
\end{enumerate}
\end{proposicion}

\addtocounter{definicion}{1}%

\section{Demostración por contrarrecíproco}

\begin{definicion}[Demostración por contrarrecíproco]\label{def:contrarreciproco}
Para demostrar $\forall x\,(P(x)\Rightarrow Q(x))$ por \emph{contrarrecíproco}, se demuestra de forma directa $\forall x\,(\neg Q(x)\Rightarrow\neg P(x))$. Es legítimo porque ambos enunciados son lógicamente equivalentes (\cref{prop:leyes}(v)).
\end{definicion}

La técnica es útil cuando la negación de la tesis aporta información más manejable que la hipótesis. Un indicio típico es que la tesis afirme algo sobre un número ($n$ es par) y la hipótesis, algo sobre una expresión de ese número ($n^{2}$ es par)~\cite[§5.1]{hammack}.

\begin{proposicion}\label{prop:cuadrado-par}
Sea $n\in\Z$. Si $n^{2}$ es par, entonces $n$ es par.
\end{proposicion}

Un intento directo habría partido de $n^{2}=2k$ y necesitaría extraer una raíz cuadrada, que no conduce a ninguna parte. El contrarrecíproco convierte el problema en el cálculo ya hecho.

\addtocounter{definicion}{1}%

\section{Congruencias}

\begin{definicion}[Congruencia]\label{def:congruencia}
Sea $n\in\N$. Dos enteros $a$ y $b$ son \emph{congruentes módulo $n$}, y se escribe $a\equiv b\pmod n$, si $n\divides(a-b)$.
\end{definicion}

Por la división euclídea, todo entero es congruente módulo $n$ con exactamente uno de $0,1,\dots,n-1$, su resto. Por ejemplo, $n$ es par si y solo si $n\equiv 0\pmod 2$.

\begin{proposicion}\label{prop:congruencias}
Sea $n\in\N$ y sean $a,b,c,d\in\Z$ con $a\equiv b\pmod n$ y $c\equiv d\pmod n$. Entonces
\begin{enumerate}
  \item $a+c\equiv b+d\pmod n$;
  \item $ac\equiv bd\pmod n$;
  \item $a^{k}\equiv b^{k}\pmod n$ para todo $k\in\N$.
\end{enumerate}
\end{proposicion}

\addtocounter{definicion}{1}%

\section{Bicondicionales y cadenas de equivalencias}

\begin{definicion}[Demostración de un bicondicional]
Para demostrar $P\Leftrightarrow Q$ se demuestran por separado $P\Rightarrow Q$ y $Q\Rightarrow P$, cada una con la técnica más conveniente (\cref{prop:leyes}(vii)). Para demostrar que varias afirmaciones $P_1,\dots,P_k$ son equivalentes basta una cadena cíclica $P_1\Rightarrow P_2\Rightarrow\dots\Rightarrow P_k\Rightarrow P_1$~\cite[§7.2]{hammack}.
\end{definicion}

\begin{proposicion}\label{prop:par-equiv}
Para $n\in\Z$, son equivalentes: (a) $n$ es par; (b) $n^{2}$ es par; (c) $n^{2}+2n$ es múltiplo de~4.
\end{proposicion}

\addtocounter{definicion}{1}%

\addtocounter{definicion}{1}%


% =====================================================================
\chapter{Reducción al absurdo y demostración por casos}
\label{tema:4}
% =====================================================================

Algunas afirmaciones no tienen la forma de un condicional, o la tienen pero ni la hipótesis ni la negación de la tesis dan un punto de partida útil. Para ellas se dispone de la reducción al absurdo: suponer lo contrario de lo que se quiere probar y llegar a una contradicción. Este tema la presenta con dos de los teoremas más antiguos de las matemáticas, la irracionalidad de $\sqrt2$ y la infinitud de los números primos, y estudia también la demostración por casos. Se siguen el capítulo~6 y las §4.4–4.5 de~\cite{hammack}, las §1.7–1.8 de~\cite{mcs}, la §3.5 de~\cite{velleman} y el primer \emph{problem set} de CS103~\cite{cs103}.

\section{Reducción al absurdo}

\begin{definicion}[Reducción al absurdo]\label{def:absurdo}
Para demostrar una proposición $P$ por \emph{reducción al absurdo} (o por contradicción) se supone $\neg P$ y se deduce una contradicción, es decir, una proposición de la forma $C\wedge\neg C$. Se concluye $P$.
\end{definicion}

La justificación es que $(\neg P\Rightarrow(C\wedge\neg C))\Rightarrow P$ es una tautología: si $P$ fuera falsa, $\neg P$ sería verdadera, y el condicional $\neg P\Rightarrow(C\wedge\neg C)$, con hipótesis verdadera y tesis falsa, sería falso. Para un enunciado condicional $P\Rightarrow Q$, la negación es $P\wedge\neg Q$ (\cref{prop:leyes}(vi)): se suponen a la vez la hipótesis y la negación de la tesis~\cite[§6.2]{hammack}.

\addtocounter{definicion}{1}%

\addtocounter{definicion}{1}%

\section{La irracionalidad de \texorpdfstring{$\sqrt2$}{raíz de 2}}

Un número real es \emph{racional} si se escribe como $a/b$ con $a,b\in\Z$ y $b\neq 0$, e \emph{irracional} en caso contrario. Todo racional puede escribirse en \emph{forma irreducible}, $a/b$ con $b>0$ y sin divisores comunes mayores que~1: basta dividir $a$ y $b$ por su máximo común divisor. (Alternativamente, se toma la representación con $b$ positivo mínimo, que existe por el principio del buen orden del \cref{tema:5}.)

\begin{teorema}\label{teo:raiz2}
$\sqrt2$ es irracional.
\end{teorema}

La demostración sigue la del teorema~1.8.1 de~\cite{mcs}. Su estructura se repite en los \cref{eje:4-raiz3,eje:4-raiz-cubica,eje:4-log}: suponer una fracción irreducible, obtener una ecuación entre enteros y encontrar un divisor común.

\section{La infinitud de los números primos}

\begin{lema}\label{lem:divisor-primo}
Todo entero $n\ge 2$ tiene al menos un divisor primo.
\end{lema}

\begin{teorema}[Euclides]\label{teo:euclides}
Existen infinitos números primos.
\end{teorema}

\begin{observacion}
La prueba \emph{no} afirma que $N$ sea primo. Con los seis primeros primos, $N=2\cdot 3\cdot 5\cdot 7\cdot 11\cdot 13+1=30\,031=59\cdot 509$ es compuesto; lo que la prueba garantiza es que sus factores primos no están en la lista (\cref{eje:4-euclides-n}).
\end{observacion}

\section{Demostración por casos}

\begin{definicion}[Demostración por casos]\label{def:casos}
Para demostrar $Q$ a partir de $P_1\vee P_2\vee\dots\vee P_k$ se demuestran por separado $P_1\Rightarrow Q$, \dots, $P_k\Rightarrow Q$. Es válido porque $(P_1\vee\dots\vee P_k)\Rightarrow Q\eqlog(P_1\Rightarrow Q)\wedge\dots\wedge(P_k\Rightarrow Q)$, generalización del \cref{eje:1-equivalencias}(c).
\end{definicion}

Lo esencial es que los casos sean \emph{exhaustivos}: su disyunción debe ser verdadera. Las divisiones más habituales son «$n$ par o impar», los restos módulo $m$ y el signo de un número real~\cite[§4.4]{hammack}.

\begin{proposicion}\label{prop:consecutivos}
Para todo $n\in\Z$, el número $n^{2}+n$ es par.
\end{proposicion}

\begin{proposicion}\label{prop:cuadrado-mod4}
Para todo $n\in\Z$, $n^{2}\equiv 0\pmod 4$ o $n^{2}\equiv 1\pmod 4$.
\end{proposicion}

\begin{teorema}[Desigualdad triangular]\label{teo:triangular}
Para todos $x,y\in\R$, $\abs{x+y}\le\abs{x}+\abs{y}$.
\end{teorema}

Cuando dos casos se tratan con el mismo argumento intercambiando los papeles de las variables, se escribe «sin pérdida de generalidad, supóngase $x\le y$» y se trata solo uno~\cite[§4.5]{hammack}. Debe usarse solo cuando la simetría es real.

\addtocounter{definicion}{1}%

\addtocounter{definicion}{1}%


% =====================================================================
\chapter{Inducción simple y fuerte; buen orden}
\label{tema:5}
% =====================================================================

Muchos enunciados del curso dependen de un número natural: la longitud de una secuencia, el número de capas de una red, el paso de un algoritmo iterativo. La inducción matemática permite demostrarlos para todos los valores a la vez, y es la técnica de demostración más usada en informática. Este tema presenta la inducción simple, la inducción fuerte y el principio del buen orden, y muestra que son tres formas de la misma idea. Se siguen el capítulo~10 de~\cite{hammack}, los capítulos~2 y~5 de~\cite{mcs} y las §6.1, 6.2 y 6.4 de~\cite{velleman}.

\section{El principio de inducción}

\begin{teorema}[Principio de inducción]\label{teo:induccion}
Sea $P(n)$ un predicado sobre $\N_0$. Si
\begin{enumerate}
  \item (\emph{caso base}) $P(0)$ es verdadera, y
  \item (\emph{paso de inducción}) para todo $n\in\N_0$, $P(n)\Rightarrow P(n+1)$,
\end{enumerate}
entonces $P(n)$ es verdadera para todo $n\in\N_0$.
\end{teorema}

La imagen habitual es la de una fila infinita de fichas de dominó: la primera cae, y cada una que cae derriba a la siguiente. En el paso de inducción, la suposición $P(n)$ se llama \emph{hipótesis de inducción}; no se está suponiendo lo que se quiere demostrar, sino probando un condicional para un $n$ arbitrario. Si el caso base es $P(n_0)$ con $n_0\in\Z$, la conclusión vale para todo $n\ge n_0$. El \cref{teo:induccion} se toma como axioma de los números naturales o se deduce del principio del buen orden (\cref{prop:bo-induccion}).

\addtocounter{definicion}{1}%

\addtocounter{definicion}{1}%

\addtocounter{definicion}{1}%

En las desigualdades, el paso de inducción suele necesitar una estimación intermedia, como aquí $2n+2\ge n+2$. Conviene escribir primero qué hay que obtener, $P(n+1)$, y buscar cómo llegar desde $P(n)$.

\addtocounter{definicion}{1}%

\section{Inducción fuerte}

\begin{teorema}[Inducción fuerte]\label{teo:induccion-fuerte}
Sea $P(n)$ un predicado sobre los enteros $n\ge n_0$. Si $P(n_0)$ es verdadera y, para todo $n\ge n_0$,
\[
\bigl(P(n_0)\wedge P(n_0+1)\wedge\dots\wedge P(n)\bigr)\Rightarrow P(n+1),
\]
entonces $P(n)$ es verdadera para todo $n\ge n_0$.
\end{teorema}

La inducción fuerte es útil cuando $P(n+1)$ depende no de $P(n)$, sino de casos anteriores más lejanos, como al descomponer un número en factores o en las recurrencias de varios términos.

\begin{teorema}[Existencia de la factorización en primos]\label{teo:factorizacion}
Todo entero $n\ge 2$ es primo o producto de primos.
\end{teorema}

\addtocounter{definicion}{1}%

\section{El principio del buen orden}

\begin{teorema}[Principio del buen orden]\label{teo:buen-orden}
Todo subconjunto no vacío de $\N_0$ tiene un elemento mínimo.
\end{teorema}

\begin{proposicion}\label{prop:bo-induccion}
El principio del buen orden implica el principio de inducción.
\end{proposicion}

La demostración anterior es el modelo de las \emph{pruebas por el menor contraejemplo}~\cite[§2.2]{mcs}, \cite[§10.3]{hammack}: para demostrar $\forall n\,P(n)$, se supone que el conjunto de contraejemplos no es vacío, se toma el menor y se construye un contraejemplo aún menor, o se llega a otra contradicción.

\begin{proposicion}[Existencia en la división euclídea]\label{prop:division-existencia}
Sean $a\in\Z$ y $b\in\N$. Existen $q,r\in\Z$ con $a=qb+r$ y $0\le r<b$.
\end{proposicion}

\section{Definiciones recursivas}

Los sumatorios, los productorios, las potencias y muchas sucesiones se definen \emph{recursivamente}: se da el primer valor y una regla para obtener cada valor a partir de los anteriores. Por ejemplo, $\sum_{k=1}^{0}a_k=0$ y $\sum_{k=1}^{n+1}a_k=\sum_{k=1}^{n}a_k+a_{n+1}$; la sucesión de Fibonacci se define por $F_1=F_2=1$ y $F_{n+2}=F_{n+1}+F_n$. La inducción es la forma natural de demostrar propiedades de objetos definidos así: el caso base comprueba la definición inicial y el paso usa la regla recursiva~\cite[§10.5]{hammack}.

\addtocounter{definicion}{1}%

\addtocounter{definicion}{1}%


% =====================================================================
\chapter{Existencia, unicidad, contraejemplos y redacción}
\label{tema:6}
% =====================================================================

Los temas anteriores se han centrado en enunciados universales. Este tema trata los enunciados existenciales y de unicidad, explica cómo se refuta un enunciado falso y termina con unas pautas para redactar demostraciones claras, que resumen todo el módulo. Se siguen los capítulos~7 y~9 y la §5.3 de~\cite{hammack}, la §3.6 de~\cite{velleman}, la §1.9 de~\cite{mcs}, la lista de comprobación de CS103~\cite{cs103} y el método de Polya~\cite{polya}.

\section{Demostraciones de existencia}

Para demostrar $\exists x\,P(x)$ basta encontrar \emph{un} elemento $a$, llamado \emph{testigo}, y comprobar $P(a)$. Una demostración de este tipo es \emph{constructiva}. A veces se prueba la existencia sin exhibir el testigo, por ejemplo por reducción al absurdo o examinando casos sin saber cuál se da; la demostración es entonces \emph{no constructiva}~\cite[§7.3–7.4]{hammack}.

\addtocounter{definicion}{1}%

\begin{proposicion}[Existencia no constructiva]\label{prop:irracional-potencia}
Existen números irracionales $a$ y $b$ tales que $a^{b}$ es racional.
\end{proposicion}

La prueba no dice cuál de los dos pares sirve. (Se sabe, por un teorema profundo de Gelfond y Schneider, que $c$ es irracional, pero la demostración anterior no lo necesita.)

\section{Demostraciones de unicidad}

Para demostrar $\exists!\,x\,P(x)$ se prueban dos cosas (\cref{def:unicidad}): la \emph{existencia}, $\exists x\,P(x)$, y la \emph{unicidad}, que se redacta así: «sean $x$ e $y$ tales que $P(x)$ y $P(y)$; entonces $x=y$». Una alternativa equivalente es suponer $x\neq y$ y llegar a una contradicción~\cite[§3.6]{velleman}.

\begin{proposicion}[Unicidad en la división euclídea]\label{prop:division-unicidad}
Sean $a\in\Z$ y $b\in\N$. Si $a=qb+r=q'b+r'$ con $q,q',r,r'\in\Z$ y $0\le r,r'<b$, entonces $q=q'$ y $r=r'$.
\end{proposicion}

Junto con la \cref{prop:division-existencia}, esto completa la demostración del \cref{teo:division}.

\addtocounter{definicion}{1}%

\section{Contraejemplos y refutaciones}

Para \emph{refutar} un enunciado se demuestra su negación~\cite[cap.~9]{hammack}. Según la forma del enunciado:
\begin{itemize}[leftmargin=1.4em]
  \item $\forall x\,P(x)$ se refuta con un \emph{contraejemplo}: un $a$ concreto con $\neg P(a)$.
  \item $\forall x\,(P(x)\Rightarrow Q(x))$ se refuta con un $a$ que cumple $P(a)$ y no cumple $Q(a)$.
  \item $\exists x\,P(x)$ se refuta demostrando $\forall x\,\neg P(x)$, lo que exige una demostración general, no un ejemplo.
\end{itemize}

\addtocounter{definicion}{1}%

\addtocounter{definicion}{1}%

\section{Cómo redactar una demostración}

Una demostración es un texto dirigido a un lector, no una lista de cálculos. Las pautas siguientes resumen las de la §5.3 de~\cite{hammack}, la §1.9 de~\cite{mcs} y la lista de comprobación de CS103~\cite{cs103}.

\begin{enumerate}[label=\arabic*.]
  \item \emph{Anunciar el plan.} Indicar qué se supone, qué se quiere demostrar y con qué técnica («Se procede por contrarrecíproco», «Por inducción sobre $n$»).
  \item \emph{Introducir cada variable.} Toda letra debe presentarse antes de usarse, con su universo: «Sea $n$ un entero impar»; «existe $k\in\Z$ con\dots». No reutilizar una letra para objetos distintos.
  \item \emph{Usar las definiciones, no repetirlas.} Desplegar la definición en el punto donde se necesita, aplicada a los objetos concretos.
  \item \emph{Escribir frases completas.} Las fórmulas forman parte de las frases. Los símbolos $\forall$, $\exists$, $\Rightarrow$ son para enunciados formales, no abreviaturas en la prosa.
  \item \emph{Que cada frase aporte.} Eliminar lo que no se usa después; justificar cada paso no evidente citando el resultado empleado.
  \item \emph{Terminar con la conclusión.} Decir explícitamente que se ha demostrado lo que se buscaba y cerrar con $\blacksquare$.
\end{enumerate}

Para \emph{encontrar} la demostración, antes de redactarla, sirve el método de cuatro fases de Polya~\cite{polya}: comprender el problema (datos, incógnita, condiciones; escribir el enunciado con cuantificadores), concebir un plan (¿se conoce un problema relacionado?, ¿sirve un caso particular o el contrarrecíproco?), ejecutarlo comprobando cada paso, y examinar la solución obtenida (¿se usaron todas las hipótesis?, ¿hay un argumento más directo?).

\addtocounter{definicion}{1}%

\addtocounter{definicion}{1}%

\addtocounter{definicion}{1}%



\backmatter
% =====================================================================
\begin{thebibliography}{99}
\addcontentsline{toc}{chapter}{Bibliografía}
\markboth{Bibliografía}{}
\small

\bibitem{hammack}
R.~Hammack.
\newblock \emph{Book of Proof}, 3.ª ed. (edición~3.4).
\newblock Richard Hammack, 2018. Libro de texto de acceso libre.
\newblock \url{https://richardhammack.github.io/BookOfProof/}

\bibitem{mcs}
E.~Lehman, F.~T. Leighton y A.~R. Meyer.
\newblock \emph{Mathematics for Computer Science}, revisión del 18 de mayo de 2015.
\newblock Texto del curso 6.042J del MIT.
\newblock \url{https://ocw.mit.edu/courses/6-042j-mathematics-for-computer-science-spring-2015/}

\bibitem{mit6042}
A.~R. Meyer y A.~Chlipala.
\newblock 6.042J: Mathematics for Computer Science.
\newblock Massachusetts Institute of Technology, MIT OpenCourseWare, primavera de 2015. Unidad~1, «Proofs».
\newblock \url{https://ocw.mit.edu/courses/6-042j-mathematics-for-computer-science-spring-2015/}

\bibitem{polya}
G.~Polya.
\newblock \emph{How to Solve It: A New Aspect of Mathematical Method}.
\newblock Princeton University Press, 1945.

\bibitem{cs103}
Stanford University.
\newblock CS103: Mathematical Foundations of Computing.
\newblock Edición archivada \texttt{cs103.1272}, consultada en octubre de 2026: clases «Mathematical Proofs», «Indirect Proofs», «Propositional Logic» y «First-Order Logic»; guías «Guide to Proofs», «Guide to Negation» y «Proofwriting Checklist»; \emph{Problem Set}~1.
\newblock \url{https://web.stanford.edu/class/cs103/}

\bibitem{velleman}
D.~J. Velleman.
\newblock \emph{How to Prove It: A Structured Approach}, 3.ª ed.
\newblock Cambridge University Press, 2019.

\end{thebibliography}


\end{document}
